\subsection{Query Delegation in Discovery Settings}
\label{sec:authority}

Dynamically integrating conjunctive sub-queries during federated query processing 
in discovery-based querying at runtime is difficult due to the unbounded nature of
the query execution. To show this, we first introduce a sufficient
condition for correctly delegating conjunctive queries in the traditional federated case.
Then we will show how deciding whether we can correctly delegate these queries in the
discovery-based querying case is impossible when not all sources have been discovered.
 
First we define our discovery-based federated query setting in 
Definition~\ref{def:setting} and~\ref{def:reach}.

\begin{definition}[Discovery setting]
\label{def:setting}
A \emph{discovery setting} is a tuple
$W = \langle \Du, \dataf, \Dd_0, \disc \rangle$: a universe of sources $\Du$, a map
$\dataf : \Du \to 2^{\Tt}$ giving the graph each one holds, a set of seeds
$\Dd_0 \subseteq \Du$, and a \emph{discovery function} $\disc$ assigning to every BGP $B$
a map $\disc_B : \Du \times 2^{\Tt} \to 2^{\Du}$, where $\disc_B(d, \dataf(d))$ are the
sources that a retrieved $d$ exposes while $B$ is evaluated.
\end{definition}

\begin{definition}[Reachable set, reference answer]
\label{def:reach}
Iterating discovery from the seeds, let
$\Dd_{i+1} := \Dd_i \cup \bigcup_{d \in \Dd_i} \disc_B\bigl(d, \dataf(d)\bigr)$. The
\emph{reachable set} is $\reach_B(W) := \bigcup_{i \geq 0} \Dd_i$, and the
\emph{reference answer} to $B$ is $\ev{B}{G_{\reach_B(W)}}$. Correctness of a delegation
is always measured against this reference answer, never against the universe $\Du$:
sources that are never reached cannot contribute to it.
\end{definition}

\noindent

Given this definition, we recover the traditional federated query processing setting,
in which all sources are known beforehand, by taking $\Dd_0 = \Du$ 
and $\disc_B \equiv \emptyset$.

Next we define an operator to determine the \emph{relevance} of a source for a query
execution setting. Relevance records, for each triple pattern and each source, whether
that source produces matches. An example of such an operator is the usage of \textsf{ASK}
queries to determine exclusive groups in FedX~\cite{schwarte2011fedx}.

\begin{definition}[Source relevance]
\label{def:relevance}
Fix a setting $W$ and a BGP $B$, and write $\Dd := \reach_B(W)$ for its reachable set;
all that follows is relative to $\Dd$. We write $\ev{t}{d}$ for the evaluation of a
triple pattern $t$ over the data of source $d$. The \emph{relevance indicator}
$\Delta : \mathcal{P} \times \Dd \to \{0,1\}$ is
\[
  \Delta(t, d) \;=\;
  \begin{cases}
    1, & \text{if } \llbracket t \rrbracket_{d} \neq \emptyset, \\[2pt]
    0, & \text{otherwise,}
  \end{cases}
\]
and $\Dd_t := \{\, d \in \Dd \mid \Delta(t,d) = 1 \,\}$ denotes the
\emph{relevant sources} of $t$.
\end{definition}

Thus, relevance means that a data source returns answers for a given triple pattern.
Next, we introduce a definition of \emph{correct} delegation of conjunctive queries to a
set of sources.

\begin{definition}[Delegation]
\label{def:delegation}
Let $B$ be a BGP, $B' \subseteq B$ the delegated fragment, and $S \subseteq \Dd$ the
(only) sources to which $B'$ is sent as a single conjunctive request, where as in
Definition~\ref{def:relevance}, $\Dd = \reach_B(W)$. Write
$G_A := \bigcup_{d \in A} \dataf(d)$ for $A \subseteq \Dd$. The delegated plan obtains for $B'$
\[
  R \;:=\; \Big( \bigcup_{d \in S} \ev{B'}{\dataf(d)} \Big)
           \;\cup\; \ev{B'}{G_{\Dd \setminus S}},
\]
and the delegation is \emph{correct} if
\[
  R \;\bowtie\; \ev{B \setminus B'}{G_{\Dd}} \;=\; \ev{B}{G_{\Dd}} .
\]
\end{definition}

This definition follows some of the definitions of federation of FedQPL~\cite{cheng2020fedqpl}.
Next we give a sufficient condition for correctly delegating conjunctive (sub)-queries to
federated sources.

\begin{theorem}[Correct delegation]
  \label{th:correct-delegation}
    Let $W$ be a \emph{federated} setting, that is $\disc_B \equiv \emptyset$, so that
    $\Dd = \reach_B(W) = \Dd_0$ is known in advance.
    Let $t_{1}, \dots, t_{n}$ be a set of triple patterns forming a BGP $B$, 
    $B' \subseteq B$ a sub-BGP, and let $d \in \Dd$. 
    For $d$ we define $\Dd^{\prime} = \Dd \setminus \{d\}$
    the set containing all other sources. 
    Then we can correctly delegate $B^{\prime}$ to $d$ if
    $\bigcup_{t_i \in B^{\prime}}
    \{\, d^{\prime} \in \Dd^{\prime} \mid \Delta(t_i, d^{\prime}) = 1 \,\} = \emptyset$
\end{theorem}
\begin{proof}
We prove the special case of Definition~\ref{def:delegation} with $S = \{d\}$, so that
$\Dd \setminus S = \Dd^{\prime}$, and we assume $B^{\prime} \neq \emptyset$.
%% This proof previously also assumed here "that distinct sources use disjoint sets of
%% blank nodes". That clause now relies on the preliminaries instead: Assumption 1
%% (Skolemised sources) excludes blank nodes from source graphs altogether, so there are
%% none left to keep disjoint, and \eqref{eq:distrib} below holds as a literal equality
%% rather than only up to blank-node renaming. Restore the clause here if that
%% assumption is ever dropped.

The argument rests on the fact that a \emph{single} triple pattern distributes over a
union of sources: for every $t$ and every $A \subseteq \Dd$,
\begin{equation}
  \ev{t}{G_{A}} \;=\; \bigcup_{e \in A} \ev{t}{\dataf(e)},
  \label{eq:distrib}
\end{equation}
since matching one pattern involves no join. Note that \eqref{eq:distrib} fails for BGPs,
which is exactly what makes delegation non-trivial.

By hypothesis no source other than $d$ is relevant to any pattern of $B^{\prime}$, so
Definition~\ref{def:relevance} gives
\begin{equation}
  \ev{t^{\prime}}{\dataf(d^{\prime})} = \emptyset
  \quad \text{for all } t^{\prime} \in B^{\prime},\ d^{\prime} \in \Dd^{\prime}.
  \label{eq:empty}
\end{equation}

\emph{Step 1.} A BGP is the join of its patterns, and each pattern distributes over the
sources by \eqref{eq:distrib}:
\[
  \ev{B}{G_{\Dd}}
  \;=\; \bigjoin_{t \in B} \ev{t}{G_{\Dd}}
  \;=\; \bigjoin_{t \in B} \ \bigcup_{e \in \Dd} \ev{t}{\dataf(e)} .
\]

\emph{Step 2.} Since $\bowtie$ is associative and commutative, the delegated fragment
splits off:
\[
  \ev{B}{G_{\Dd}} \;=\;
  \Bigl( \bigjoin_{t^{\prime} \in B^{\prime}} \bigcup_{e \in \Dd}
         \ev{t^{\prime}}{\dataf(e)} \Bigr)
  \bowtie \ev{B \setminus B^{\prime}}{G_{\Dd}} .
\]

\emph{Step 3.} As $\Dd = \Dd^{\prime} \uplus \{d\}$, the inner union splits, and by
\eqref{eq:empty} its $\Dd^{\prime}$ part is empty:
\[
  \bigcup_{e \in \Dd} \ev{t^{\prime}}{\dataf(e)}
  \;=\; \underbrace{\bigcup_{d^{\prime} \in \Dd^{\prime}}
        \ev{t^{\prime}}{\dataf(d^{\prime})}}_{=\,\emptyset}
        \;\cup\; \ev{t^{\prime}}{\dataf(d)}
  \;=\; \ev{t^{\prime}}{\dataf(d)} .
\]
Substituting into Step 2 and reassembling the join,
\begin{equation}
  \ev{B}{G_{\Dd}} \;=\;
  \ev{B^{\prime}}{\dataf(d)} \bowtie \ev{B \setminus B^{\prime}}{G_{\Dd}} .
  \label{eq:ref}
\end{equation}

\emph{Step 4.} It remains to identify $R$ from Definition~\ref{def:delegation}. Its
second term vanishes, again by \eqref{eq:distrib} and \eqref{eq:empty} and using
$B^{\prime} \neq \emptyset$:
\[
  \ev{B^{\prime}}{G_{\Dd \setminus S}}
  \;=\; \bigjoin_{t^{\prime} \in B^{\prime}} \ \bigcup_{d^{\prime} \in \Dd^{\prime}}
        \ev{t^{\prime}}{\dataf(d^{\prime})}
  \;=\; \bigjoin_{t^{\prime} \in B^{\prime}} \emptyset
  \;=\; \emptyset .
\]
With $S = \{d\}$ the first term of $R$ is the single request to $d$, so
$R = \ev{B^{\prime}}{\dataf(d)} \cup \emptyset = \ev{B^{\prime}}{\dataf(d)}$.

\emph{Conclusion.} Substituting $R$ into \eqref{eq:ref},
\[
  R \bowtie \ev{B \setminus B^{\prime}}{G_{\Dd}} \;=\; \ev{B}{G_{\Dd}},
\]
which is the correctness condition of Definition~\ref{def:delegation}.
\end{proof}

Theorem~\ref{th:correct-delegation} gives us a sufficient condition for delegating 
conjuctive sub-queries in the case where the reachable set $\Dd$ is known in advance. However, in discovery-based
querying the sources are discovered at runtime, goverened by some discovery process which returns
during the query execution. As such, determining if we can correctly delegate depends on the
\emph{state} of the discovery process, which is defined in Definition~\ref{def:discovery-state}.

\begin{definition}[Discovery state]
  \label{def:discovery-state}
  The state of a discovery setting $W = \langle \Du, \dataf, \Dd_0, \disc \rangle$ at
  time $T$ is the tuple
  \[
    W(T) \;=\; \langle \Du, \dataf, \Dd_0, \disc, \Dd(T) \rangle,
  \]
  the setting $W$ augmented with $\Dd(T)$, the set of sources that $\dataf$ has been
  called on by time $T$, subject to
  \[
    \Dd_0 \;\subseteq\; \Dd(T) \;\subseteq\; \Dd(T') \;\subseteq\; \reach_B(W)
    \qquad \text{for all } T \leq T',
  \]
  i.e.\ the seeds are discovered from the start, discovery only ever adds sources, and
  it never adds one outside what $\disc$ can reach. For any given $W(T)$ the discovery
  \emph{frontier} is
  \[
    F(T) \;=\; \Bigl( \bigcup_{d_f \in \Dd(T)} \disc\bigl(d_f, \dataf(d_f)\bigr) \Bigr)
                \setminus \Dd(T),
  \]
  the sources that are discovered but do not yet have $\dataf$ called on them.
  Discovery is complete when $F(T) = \emptyset$.
\end{definition}

Next, we introduce the concept of an observation of an engine of the discovery state 
in Definition~\ref{def:obs-equiv}, which
denotes what the engine actually can \emph{know} about its discovery setting at given time $T$.

\begin{definition}[Observationally equivalent states]
  \label{def:obs-equiv}
  A discovery state $W(T)$ records the whole world, but a query engine at time $T$ has
  access only to the \emph{observation}
  \[
    O(T) \;=\; \bigl\langle\, \Dd(T),\; \dataf|_{\Dd(T)} \,\bigr\rangle,
  \]
  the sources it has retrieved together with their contents. 
  Two states $W_1(T)$ and $W_2(T)$ are \emph{observationally equivalent}, written
  $W_1(T) \sim W_2(T)$, when $O_1(T) = O_2(T)$. Note that $F(T)$ is computable from
  $O(T)$ and $\disc$: the extent of the frontier is observable, its contents are not. A
  property of states is \emph{determined at $T$} when all states observationally
  equivalent to $W(T)$ agree on it; an engine at time $T$ can establish a property only
  if it is determined at $T$.
\end{definition}

\begin{theorem}[Correct delegation is undetermined under incomplete discovery]
  \label{th:undetermined}
  Let $W(T)$ be a discovery state for which $F(T) \neq \emptyset$. Then there exist a
  BGP $B$, a sub-BGP $B^{\prime} \subseteq B$, a set of sources
  $S \subseteq \Dd(T)$, and two states $W_1(T)$ and $W_2(T)$, both observationally
  equivalent to $W(T)$, such that $B^{\prime}$ is correctly delegated to $S$ in $W_1(T)$
  but not in $W_2(T)$. Since the engine's input is $O(T)$ and all three states share it,
  no procedure can decide whether the delegation is correct.
\end{theorem}

\begin{proof}
  We construct two states $W_1(T)$ and $W_2(T)$, both observationally equivalent to
  $W(T)$, in which $B^{\prime}$ is delegated to $S$ correctly and incorrectly
  respectively. The two agree on $\Dd_0$, $\disc$ and $\Dd(T)$, and on $\dataf$ over
  $\Dd(T)$; they differ only beyond the frontier, which is non-empty by hypothesis.

  Take $\Dd(T) = S = \{d_1\}$ and
  \[
    B^{\prime} \;=\; B \;=\; \bigl\{\, (?s, p_1, ?o_1),\ (?s, p_2, ?o_2) \,\bigr\},
  \]
  \[
    \dataf(d_1) \;=\; \bigl\{\, (s_1, p_1, o_1),\ (s_1, p_2, o_2),\ (s_2, p_1, o_3) \,\bigr\}.
  \]
  Fix some $d^{*} \in F(T)$, and let every reachable source other than $d_1$ and $d^{*}$
  hold no data. The two states are then obtained by setting
  \[
    \dataf(d^{*}) \;=\;
    \begin{cases}
      \bigl\{\, (s_3, p_2, o_4) \,\bigr\} & \text{in } W_1(T), \\[2pt]
      \bigl\{\, (s_2, p_2, o_4) \,\bigr\} & \text{in } W_2(T).
    \end{cases}
  \]
  Since $d^{*} \notin \Dd(T)$, both states induce the same observation, and so
  $W_1(T) \sim W_2(T) \sim W(T)$ holds. The reachable set differs between the two, so we
  write $G_i$ for $G_{\reach_B(W_i)}$, the union of all data reachable in $W_i(T)$. As
  only $d_1$ and $d^{*}$ hold data,
  \[
    G_1 = \dataf(d_1) \cup \bigl\{\, (s_3, p_2, o_4) \,\bigr\},
    \qquad
    G_2 = \dataf(d_1) \cup \bigl\{\, (s_2, p_2, o_4) \,\bigr\}.
  \]

  Write
  \begin{align*}
    \mu_1 &= \{?s \mapsto s_1,\ ?o_1 \mapsto o_1,\ ?o_2 \mapsto o_2\}, \\
    \mu_2 &= \{?s \mapsto s_2,\ ?o_1 \mapsto o_3,\ ?o_2 \mapsto o_4\}.
  \end{align*}
  Neither added triple matches $(?s, p_1, ?o_1)$, so
  $\ev{B^{\prime}}{\dataf(d^{*})} = \emptyset$ in both states. By
  Definition~\ref{def:delegation} we therefore have
  \[
    R_1 \;=\; R_2 \;=\; \ev{B^{\prime}}{\dataf(d_1)} \;=\; \{\mu_1\}.
  \]
  Since $B \setminus B^{\prime} = \emptyset$ we have
  $\ev{B \setminus B^{\prime}}{G_i} = \{\mu_\emptyset\}$, the join identity, so the
  correctness condition reduces to $R_i = \ev{B}{G_i}$.

  In $W_1(T)$ the added triple has subject $s_3$, which no $p_1$-triple mentions, so it
  contributes no solution and
  \[
    \ev{B}{G_1} \;=\; \{\mu_1\} \;=\; R_1 :
  \]
  the delegation is correct. In $W_2(T)$ the added triple joins with
  $(s_2, p_1, o_3) \in \dataf(d_1)$ on $?s$, so
  \[
    \ev{B}{G_2} \;=\; \{\mu_1, \mu_2\} \;\neq\; R_2 :
  \]
  the delegation is not correct.

  The triple at $d^{*}$ yields no solution on its own, and $(s_2, p_1, o_3)$ yields none
  within $d_1$; only their join does. A single conjunctive request to $d_1$ cannot produce
  it, and neither can one to $d^{*}$, so the solution is lost precisely because the join
  crosses the frontier. As $W_1(T)$ and $W_2(T)$ present the engine with the same
  observation $O(T)$ while disagreeing on correctness, no procedure reading only $O(T)$
  can decide it.
\end{proof}


\begin{lemma}[An empty frontier certifies complete discovery]
  \label{lem:frontier}
  $F(T) = \emptyset$ if and only if $\Dd(T) = \reach_B(W)$.
\end{lemma}

\begin{proof}
  $(\Rightarrow)$ The inclusion $\Dd(T) \subseteq \reach_B(W)$ holds by
  Definition~\ref{def:discovery-state}. For the converse, $F(T) = \emptyset$ says exactly
  that
  \[
    \disc\bigl(d_f, \dataf(d_f)\bigr) \subseteq \Dd(T)
    \qquad \text{for every } d_f \in \Dd(T),
  \]
  i.e.\ that $\Dd(T)$ is closed under $\disc$; and $\Dd_0 \subseteq \Dd(T)$,
  again by Definition~\ref{def:discovery-state}. By induction, if
  $\Dd_i \subseteq \Dd(T)$ then
  \[
    \Dd_{i+1} \;=\; \Dd_i \cup \bigcup_{d \in \Dd_i} \disc\bigl(d, \dataf(d)\bigr)
    \;\subseteq\; \Dd(T)
  \]
  by closure, so $\reach_B(W) = \bigcup_{i \geq 0} \Dd_i \subseteq \Dd(T)$.

  $(\Leftarrow)$ If $d \in \reach_B(W)$ then $d \in \Dd_i$ for some $i$, hence
  \[
    \disc\bigl(d, \dataf(d)\bigr) \subseteq \Dd_{i+1} \subseteq \reach_B(W).
  \]
  Thus $\reach_B(W)$ is closed under $\disc$, and $F(T) = \emptyset$ follows.
\end{proof}

\todo{Don't like the text here}
is what makes the frontier useful: $F(T) = \emptyset$ is a
condition the engine can test, whereas $\Dd(T) = \reach_B(W)$ is a statement about a
set it cannot inspect. Together with Theorem~\ref{th:undetermined} it fixes exactly when
delegation can be established at all.

\begin{corollary}[Delegation requires complete discovery]
  \label{cor:complete-discovery}
  The correctness of every delegation of a sub-BGP $B^{\prime} \subseteq B$ to a set
  $S \subseteq \Dd(T)$ is determined at $T$ if and only if $F(T) = \emptyset$.
\end{corollary}

\begin{proof}
  $(\Leftarrow)$ Let $F(T) = \emptyset$ and let $W_1(T), W_2(T)$ be observationally
  equivalent to $W(T)$. As $F(T)$ is computable from $O(T)$, both have an empty frontier,
  so $\Dd = \reach_B(W) = \Dd(T)$ in both by Lemma~\ref{lem:frontier}. Then
  $G_{\Dd} = \bigcup_{e \in \Dd(T)} \dataf(e)$ agrees in both, and since
  $S \subseteq \Dd(T) = \Dd$, so do $G_{\Dd \setminus S}$ and $\dataf(d)$ at each
  $d \in S$. Every term of the condition of Definition~\ref{def:delegation} therefore
  agrees, and with it the verdict.

  $(\Rightarrow)$ By contraposition. If $F(T) \neq \emptyset$, then
  Theorem~\ref{th:undetermined} supplies a $B^{\prime}$, an $S$, and two states
  observationally equivalent to $W(T)$ that disagree on the correctness of delegating
  $B^{\prime}$ to $S$. That delegation is not determined at $T$.
\end{proof}

